\documentclass[a4paper,11pt]{article}
\usepackage[utf8]{inputenc}
\usepackage[T2A]{fontenc}
\usepackage[english,russian]{babel}
\usepackage[margin=2.5cm]{geometry}
\usepackage{amsmath,amssymb,amsthm}
\usepackage{stmaryrd}
\usepackage{proof}
\usepackage{listings}

\newtheorem{thm}{Теорема}[section]
\newtheorem{lmm}[thm]{Лемма}
\theoremstyle{definition}
\newtheorem{dfn}[thm]{Определение}

\newcommand{\sem}[2]{\llbracket #1 \rrbracket_{#2}}
\newcommand{\cmpl}{\mathsf{c}}
\newcommand{\proves}{\vdash}
\newcommand{\True}{\text{И}}
\newcommand{\False}{\text{Л}}

\begin{document}

\section{Общие замечания}
\subsection{О тексте}
Данный текст представляет из себя краткий конспект лекций по курсу
«Математическая логика», рассказанных студентам ИТМО (группы M3234--M3239) в
2017--2018 учебном году.

\subsection{Общие определения и обозначения}
Прежде чем приступить к изложению содержательного материала, введём несколько
базовых определений и обозначений, которые должны быть уже знакомы читателю.
\begin{enumerate}
\item Множество всех подмножеств обозначим как $\mathcal{P}$:
  $\mathcal{P}(X) = \{C \mid C \subseteq X\}$.
\item Упорядоченную пару каких-либо значений $a$ и $b$ будем обозначать как
  $\langle a, b\rangle$.
\item Пусть дано некоторое частично-упорядоченное отношением $\sqsubseteq$
  множество $S$. Наименьшим (наибольшим) элементом множества назовём такой
  элемент $t \in S$, что для любого $s \in S$ выполнено $t \sqsubseteq s$
  ($s \sqsubseteq t$).
\item Пусть дано некоторое частично-упорядоченное отношением $\sqsubseteq$
  множество $S$. Минимальным (максимальным) элементом множества назовём такой
  элемент $t \in S$, что не существует большего (меньшего) $s \in S$. Иными
  словами, нет такого $s$, что $s \sqsubseteq t$ ($t \sqsubseteq s$) и $s \ne t$.

  Заметим, что наименьшее (наибольшее) значение всегда единственное, а
  минимальных (максимальных) значений может быть много.
\end{enumerate}

\section{Общая топология}
Мы начинаем курс немного издалека: от некоторых базовых тем общей топологии. С
одной стороны, эти знания пригодятся нам дальше в курсе, с другой~--- есть
надежда, что они настроят слушателей курса на правильный лад.

\begin{dfn}
Топологическим пространством мы назовём упорядоченную пару множеств
$\langle X, \Omega\rangle$, где $\Omega \subseteq \mathcal{P}(X)$, отвечающую
следующим трём свойствам:
\begin{enumerate}
\item Какое бы ни было семейство множеств $\{A_\alpha\}$, где
  $A_\alpha \in \Omega$, выполнено $\bigcup_\alpha \{A_\alpha\} \in \Omega$.
\item Какое бы ни было конечное семейство множеств $\{A_1, \dots, A_n\}$, где
  $A_i \in \Omega$, выполнено $A_1 \cap A_2 \cap \dots \cap A_n \in \Omega$.
\item $\varnothing \in \Omega$, $X \in \Omega$.
\end{enumerate}
\end{dfn}

\begin{dfn}
Пусть дано топологическое пространство $\langle X, \Omega\rangle$. Тогда любое
множество $A \in \Omega$ назовём открытым. Если же $X \setminus A \in \Omega$,
то такое множество назовём замкнутым.
\end{dfn}

\begin{thm}
Следующие объекты являются топологическими пространствами:
\begin{enumerate}
\item Топология стрелки:
  $\langle \mathbb{R}, \{(x, +\infty) \mid x \in \{-\infty\} \cup \mathbb{R} \cup \{+\infty\}\}\rangle$.
\item Дискретная топология на множестве $X$: $\langle X, \mathcal{P}(X)\rangle$.
\item Топология Зарисского на множестве $X$:
  $\langle X, \{A \in \mathcal{P}(X) \mid (X \setminus A) \text{ --- конечно}\}\rangle$.
\end{enumerate}
\end{thm}

\begin{dfn}
Внутренностью множества $X$ (обозначается как $\operatorname{int} X$) мы назовём
наибольшее по включению открытое подмножество $X$. Замыканием множества $X$
(обозначается как $\operatorname{cl} X$) мы назовём наименьшее по включению
замкнутое надмножество $X$.
\end{dfn}

\begin{dfn}
Базой топологического пространства $\langle X, \Omega\rangle$ назовём любое такое
семейство множеств $B$, что каждое открытое множество представляется
объединением некоторого подмножества $B$. Или, в формальной записи,
$\Omega = \{\bigcup S \mid S \subseteq B\}$. Также будем говорить, что данная
база $B$ задаёт топологическое пространство $\langle X, \Omega\rangle$.
\end{dfn}

\begin{thm}
Классическая топология Евклидова пространства $\mathbb{R}$: множество
\[ B = \{(a, b) \mid a, b \in \mathbb{R}\} \]
является базой Евклидова пространства.
\end{thm}

\begin{dfn}
Топологическое пространство $\langle X, \Omega\rangle$ назовём связным, если
единственные одновременно открытые и замкнутые множества в нём~--- $\varnothing$
и $X$.
\end{dfn}

\begin{thm}
Топологическое пространство $\langle X, \Omega\rangle$ связно тогда и только
тогда, когда в нём нет двух непустых открытых множеств $A$ и $B$, что
$A \cup B = X$ и $A \cap B = \varnothing$.
\end{thm}

\begin{dfn}
Назовём частичным порядком ($\sqsubseteq$) на множестве $X$ любое рефлексивное,
транзитивное и антисимметричное отношение на нём.
\end{dfn}

\begin{dfn}
Рассмотрим множество $X$ с заданным на нём частичным порядком $\sqsubseteq$.
Рассмотрим множество
$B_\sqsubseteq = \{\{t \in X \mid x \sqsubseteq t\} \mid x \in X\}$. Тогда
топологическое пространство $X_\sqsubseteq$, задаваемое базой топологии
$B_\sqsubseteq$, мы назовём топологией частичного порядка ($\sqsubseteq$) на $X$.
\end{dfn}

\begin{thm}
При любом выборе $X$ и $\sqsubseteq$ $X_\sqsubseteq$ является топологическим
пространством.
\end{thm}

\begin{dfn}
Пусть задано топологическое пространство $\langle X, \Omega\rangle$, и пусть
задано множество $A \subseteq X$. Тогда рассмотрим
$\Omega_A = \{S \cap A \mid S \in \Omega\}$. Будем называть топологическое
пространство $\langle A, \Omega_A\rangle$ пространством с топологией,
индуцированной пространством $\langle X, \Omega\rangle$.
\end{dfn}

\begin{thm}
При любом выборе топологического пространства $\langle X, \Omega\rangle$ и $A$
(подмножества $X$) пространство с индуцированной топологией
$\langle A, \Omega_A\rangle$ является топологическим пространством.
\end{thm}

\begin{dfn}
Пусть задано топологическое пространство $\langle X, \Omega\rangle$, и пусть
$A \subseteq X$. Тогда множество $A$ называется связным, если оно связно как
пространство с индуцированной пространством $\langle X, \Omega\rangle$
топологией.
\end{dfn}

\begin{thm}
Рассмотрим ациклический граф $G$ с множеством вершин $V$. Построим по нему
отношение: положим, что $x \sqsubseteq y$, если имеется путь из $x$ в $y$ (в
частности, будем считать, что всегда есть путь из $x$ в $x$). Тогда граф слабо
связен тогда и только тогда, когда связно соответствующее топологическое
пространство частичного порядка.
\end{thm}

\section{Исчисление высказываний}
Матлогика~--- это наука о правильных математических рассуждениях, а поскольку
рассуждения обычно ведутся на каком-то языке, то она неразрывно связана с идеей
двух языков: языка исследователя (или иначе его называют метаязыком), и
предметного языка. Как следует из названий, языком исследователя пользуемся мы,
формулируя утверждения или доказывая теоремы о разных способах математических
рассуждений, или просто их обсуждая. Сами же математические рассуждения,
собственно и составляющие предмет исследования, формализованы в некотором
предметном языке.

Мы начнём с очень простого предметного языка~--- языка исчисления высказываний.
Элементами (строками) данного языка являются некоторые выражения (формулы), по
структуре очень похожие на арифметические, которые называются высказываниями.

Каждое высказывание~--- это либо пропозициональная переменная~--- большая буква
латинского алфавита, возможно, с цифровым индексом, либо оно составлено из
одного или двух высказываний меньшего размера, соединённых логической связкой.

Связок в языке мы определим 4 (хотя при необходимости этот список может быть в
любой момент изменен).
\begin{itemize}
\item отрицание: если $\alpha$~--- высказывание, то $\neg\alpha$~--- тоже
  высказывание.
\item конъюнкция: если $\alpha$ и $\beta$~--- высказывания, то $\alpha\mathbin{\&}\beta$~---
  тоже высказывание.
\item дизъюнкция: если $\alpha$ и $\beta$~--- высказывания, то $\alpha\vee\beta$~---
  тоже высказывание.
\item импликация: если $\alpha$ и $\beta$~--- высказывания, то $\alpha\to\beta$~---
  тоже высказывание.
\end{itemize}

Также в языке можно использовать скобки вокруг выражений: если $\alpha$~---
высказывание, то $(\alpha)$~--- тоже высказывание. Если из расстановки скобок не
следует иное, мы будем учитывать приоритет связок (связки в перечислении выше
указаны в порядке убывания приоритета). Также, конъюнкцию и дизъюнкцию мы будем
считать левоассоциативной (скобки в цепочке одинаковых связок расставляются
слева направо: $(A\vee B)\vee C$), а импликацию~--- правоассоциативной:
$A\to(B\to C)$.

Высказывания, подробности которых нас не интересуют, мы будем обозначать
начальными буквами греческого алфавита ($\alpha$, $\beta$, $\gamma$ и т.п.). Ещё
мы будем называть такие высказывания метапеременными. Одинаковым буквам всегда
соответствуют одинаковые высказывания, однако, разным буквам не обязаны
соответствовать разные высказывания. При подстановке выражений вместо
метапеременных мы всегда предполагаем, что эти выражения взяты в скобки.

Покажем эти правила на примере. Выражение
\[ \alpha \to \neg\beta\mathbin{\&}\gamma \vee \delta \vee \varepsilon \to \zeta \vee \iota \]
нужно воспринимать так:
\[ (\alpha) \to (((((\neg(\beta))\mathbin{\&}(\gamma)) \vee (\delta)) \vee (\varepsilon)) \to ((\zeta) \vee (\iota))) \]

\subsection{Оценка высказываний}
Процесс «вычисления» значения высказываний имеет совершенно естественное
определение. Мы фиксируем некоторое множество истинностных значений $V$, для
начала мы в качестве такого множества возьмем множество $\{\True, \False\}$,
здесь $\True$ означает истину, а $\False$~--- ложь. Всем пропозициональным
переменным мы приписываем некоторое значение, а затем рекурсивно вычисляем
значение выражения естественным для указанных связок образом.

Пусть $P$~--- множество пропозициональных переменных языка. Тогда функцию
$M : P \to V$, приписывающую истинностное значение пропозициональным
переменным, мы назовём моделью (иначе: интерпретацией или оценкой переменных).

Функцию, сопоставляющую высказыванию $\alpha$ и модели $M$ истинностное
значение, мы назовём оценкой высказывания и будем это записывать так:
$\sem{\alpha}{M}$. Обычно для указания модели $M$ мы будем перечислять значения
пропозициональных переменных, входящих в формулу:
$\sem{P \to Q}{P:=\False,\,Q:=\True} = \True$. Если конкретная модель ясна из
контекста или несущественна для изложения, мы можем указание на модель
опустить: $\llbracket P \to P \rrbracket = \True$.

Если $\sem{\alpha}{M} = \True$, то мы будем говорить, что высказывание $\alpha$
истинно в модели $M$, или, иными словами, $M$~--- модель высказывания $\alpha$.

Тавтологией или общезначимым высказыванием мы назовём высказывание, истинное в
любой модели. На языке исследователя общезначимость высказывания $\alpha$ можно
кратко записать как $\models \alpha$.

Указанный способ оценки высказываний мы будем называть классическим. В
дальнейшем мы будем брать необычные множества истинностных значений и будем
давать неожиданный смысл связкам, однако, классический способ будет всегда
подразумеваться, если не указано иного. Если же мы захотим сделать на этом
акцент, мы будем говорить о классических моделях исчисления высказываний,
подразумевая, что если мы приписываем переменным классические значения истина и
ложь, то и высказывание целиком мы оцениваем тоже по классическим правилам.

\subsection{Доказательства}
В любой теории есть некоторые утверждения (аксиомы), которые принимаются без
доказательства. В исчислении высказываний мы должны явно определить список всех
возможных аксиом. Например, мы можем взять утверждение $A\mathbin{\&}B \to A$ в
качестве аксиомы. Однако, есть множество аналогичных утверждений, например,
$B\mathbin{\&}A \to B$, которые не отличаясь по сути, отличаются по записи, и
формально говоря, являются другими утверждениями.

Для решения вопроса мы введём понятие схемы аксиом~--- некоторого обобщённого
шаблона, подставляя значения в который, мы получаем различные, но схожие
аксиомы. Например, схема аксиом $\alpha\mathbin{\&}\beta \to \alpha$ позволяет
получить как аксиому $A\mathbin{\&}B \to A$ (при подстановке $\alpha := A$,
$\beta := B$), так и аксиому $B\mathbin{\&}A \to B$.

Возьмем следующие схемы аксиом для исчисления высказываний.
\begin{align*}
&(1)\ \alpha \to \beta \to \alpha\\
&(2)\ (\alpha \to \beta) \to (\alpha \to \beta \to \gamma) \to (\alpha \to \gamma)\\
&(3)\ \alpha \to \beta \to \alpha\mathbin{\&}\beta\\
&(4)\ \alpha\mathbin{\&}\beta \to \alpha\\
&(5)\ \alpha\mathbin{\&}\beta \to \beta\\
&(6)\ \alpha \to \alpha \vee \beta\\
&(7)\ \beta \to \alpha \vee \beta\\
&(8)\ (\alpha \to \gamma) \to (\beta \to \gamma) \to (\alpha \vee \beta \to \gamma)\\
&(9)\ (\alpha \to \beta) \to (\alpha \to \neg\beta) \to \neg\alpha\\
&(10)\ \neg\neg\alpha \to \alpha
\end{align*}

Напомним, что импликация~--- правоассоциативная операция, поэтому скобки в
схеме аксиом 1, например, расставляются так: $(\alpha) \to ((\beta) \to (\alpha))$.

Помимо аксиом, нам требуется каким-то образом научиться преобразовывать одни
верные утверждения в другие. Сделаем это с помощью правил вывода. У нас пока
будет одно правило вывода~--- Modus Ponens. Это также схема, она позволяет при
доказанности двух формул $\psi$ и $\psi \to \varphi$ считать доказанной формулу
$\varphi$.

\begin{dfn}
Доказательство в исчислении высказываний~--- это некоторая конечная
последовательность выражений $\alpha_1, \alpha_2 \dots \alpha_n$ из языка $L$,
такая, что каждое из утверждений $\alpha_i$ ($1 \le i \le n$) либо является
аксиомой, либо получается из других утверждений
$\alpha_{P_1}, \alpha_{P_2} \dots \alpha_{P_k}$ ($P_1 \dots P_k < i$) по
правилу вывода.
\end{dfn}

\begin{dfn}
Высказывание $\alpha$ называется доказуемым, если существует доказательство
$\alpha_1, \alpha_2 \dots \alpha_k$, и в нем $\alpha_k$ совпадает с $\alpha$.
\end{dfn}

Вообще, схемы аксиом и правила вывода существуют для удобства задания
исчисления. В дальнейшем будет очень неудобно возиться с этими объектами.
Поэтому мы считаем, что в исчислении имеется бесконечно много аксиом и правил
вывода, которые порождаются подстановкой всех возможных формул вместо параметров
в схемы.

В качестве сокращения записи в языке исследователя мы будем писать
$\proves \alpha$, чтобы сказать, что $\alpha$ доказуемо.

Традиционно правило вывода Modus Ponens записывают так:
\[ \infer{\beta}{\alpha & \alpha \to \beta} \]

\section{Теорема о дедукции}
Соглашение об обозначениях. Будем обозначать буквами $\Gamma$, $\Delta$,
$\Sigma$, $\Pi$ списки формул (возможно, пустые).

\begin{dfn}[Вывод из допущений]
Пусть $\Gamma$~--- некоторый список высказываний, а $\alpha$~--- некоторое
высказывание. Тогда мы будем говорить, что высказывание $\alpha$ выводимо из
$\Gamma$ (и записывать это как $\Gamma \proves \alpha$), если существует такая
последовательность высказываний $\alpha_1, \alpha_2, \dots \alpha_{n-1},
\alpha$ (называемая выводом $\alpha$ из $\Gamma$), что каждое из высказываний
$\alpha_i$~--- это
\begin{itemize}
\item либо аксиома,
\item либо получается по правилу Modus Ponens из предыдущих высказываний,
\item либо высказывание из списка $\Gamma$.
\end{itemize}
Элементы $\Gamma$ мы будем называть допущениями. Также эти элементы называют
предположениями или гипотезами.
\end{dfn}

В свете данного определения можно заметить, что доказательство~--- это вывод из
пустого списка допущений.

\begin{thm}[Теорема о дедукции]
Утверждение $\Gamma \proves \alpha \to \beta$ справедливо тогда и только тогда,
когда справедливо, что $\Gamma, \alpha \proves \beta$.
\end{thm}

Для доказательства рассмотрим следующую лемму:

\begin{lmm}
$\proves \alpha \to \alpha$.
\end{lmm}
\begin{proof}
\begin{align*}
&(1)\ \alpha \to (\alpha \to \alpha) &&\text{Сх. акс. 1}\\
&(2)\ (\alpha \to (\alpha \to \alpha)) \to (\alpha \to ((\alpha \to \alpha) \to \alpha)) \to (\alpha \to \alpha) &&\text{Сх. акс. 2}\\
&(3)\ (\alpha \to ((\alpha \to \alpha) \to \alpha)) \to (\alpha \to \alpha) &&\text{M.P. 1,2}\\
&(4)\ (\alpha \to ((\alpha \to \alpha) \to \alpha)) &&\text{Сх. акс. 1}\\
&(5)\ \alpha \to \alpha &&\text{M.P. 4,3}
\end{align*}
\end{proof}

\begin{proof}[Доказательство теоремы о дедукции]
Сперва докажем прямое следствие. Для этого нам достаточно научиться по любому
выводу $\alpha \to \beta$ из $\Gamma$ строить вывод $\beta$ из $\Gamma, \alpha$.
Возьмем вывод формулы $\alpha \to \beta$, то есть некоторую последовательность
формул $\delta_1 \dots \delta_{m-1}; \alpha \to \beta$. Добавив к выводу 2
формулы, получаем требуемый вывод:
\begin{align*}
&(1)\ \delta_1\\
&\quad\dots\\
&(m-1)\ \delta_{m-1}\\
&(m)\ \alpha \to \beta\\
&(m+1)\ \alpha &&\text{«Свежедобавленная» аксиома}\\
&(m+2)\ \beta &&\text{M.P. } m, m+1
\end{align*}

Теперь докажем обратное. Нам необходимо построить вывод утверждения
$\Gamma \proves \alpha \to \beta$ по имеющемуся выводу
$\delta_1 \dots \delta_{m-1}, \beta$. Мы поступим так: сперва набросаем план
вывода~--- разместим по тексту «ключевые» формулы, которые потом дополним до
полноценного вывода промежуточными формулами.

План вывода будет такой:
\begin{align*}
&(1)\ \Gamma \proves \alpha \to \delta_1\\
&\quad\dots\\
&(m-1)\ \Gamma \proves \alpha \to \delta_{m-1}\\
&(m)\ \Gamma \proves \alpha \to \beta
\end{align*}

Теперь надо дополнить его до полноценного вывода. Будем рассматривать формулы
подряд и перед каждой формулой добавлять некоторое количество формул,
обосновывающих соответствующий шаг доказательства. Рассмотрим формулу номер $i$.
Возможны следующие варианты:
\begin{enumerate}
\item $\delta_i$~--- это аксиома или предположение, входящее в $\Gamma$. Тогда
  перед этой формулой вставим формулы $\delta_i$ и
  $\delta_i \to (\alpha \to \delta_i)$, и окажется, что $i$-я формула выводится
  из предыдущих двух формул путем применения правила Modus Ponens.
\item $\delta_i$ совпадает с $\alpha$. Тогда мы вставим перед ней 4 первые
  формулы из леммы, и $\delta_i \to \alpha$ будет получаться по правилу Modus
  Ponens.
\item $\delta_i$ выводится по правилу Modus Ponens из каких-то других
  утверждений $\delta_j$ и $\delta_k$ (при этом $\delta_k \equiv \delta_j \to
  \delta_i$), где $j < i$ и $k < i$. Покажем, что $\alpha \to \delta_i$ тоже
  может быть выведена из утверждений $\alpha \to \delta_j$ и $\alpha \to
  (\delta_j \to \delta_i)$.

  Для этого добавим два высказывания:
  \begin{align*}
  &(\alpha \to \delta_j) \to ((\alpha \to (\delta_j \to \delta_i)) \to (\alpha \to \delta_i)) &&\text{Сх. акс. 2}\\
  &((\alpha \to (\delta_j \to \delta_i)) \to (\alpha \to \delta_i)) &&\text{M.P. из } j \text{ и } i-6
  \end{align*}
\end{enumerate}
\end{proof}

По аналогии мы можем рассмотреть отношение следования. Будем говорить, что
высказывание $\alpha$ следует из высказываний $\Gamma$, если при любой оценке
пропозициональных переменных, входящих в высказывания $\Gamma$ и $\alpha$, на
которой все высказывания из $\Gamma$ истинны, $\alpha$ также истинно.
Записывать, что $\alpha$ следует из $\Gamma$, будем так: $\Gamma \models \alpha$.

\section{Теорема о полноте исчисления высказываний}
\begin{dfn}
Введем обозначение. Пусть $\alpha$~--- это некоторое высказывание, а $x$~---
некоторое истинностное значение. Тогда обозначим за $[x]\,\alpha$ высказывание
$\alpha$, если $x$~--- истина, и $\neg(\alpha)$, если $x$~--- ложь. Также, если
формула $\alpha$~--- это формула с $n$ пропозициональными переменными $P_1
\dots P_n$, и $x_1 \dots x_n$~--- некоторые истинностные значения, то за
$\sem{\alpha}{P_1:=x_1,\dots P_n:=x_n}$ обозначим значение формулы $\alpha$ при
подстановке значений $x_1 \dots x_n$ вместо переменных $P_1 \dots P_n$.
\end{dfn}

\begin{lmm}
Если $\Gamma, \Sigma \proves \alpha$, то $\Gamma, \Delta, \Sigma \proves
\alpha$. Если $\Gamma, \Delta, \Sigma, \Pi \proves \alpha$, то $\Gamma, \Sigma,
\Delta, \Pi \proves \alpha$.
\end{lmm}
\begin{proof}
Упражнение.
\end{proof}

\begin{lmm}
Если справедливы 3 утверждения: $\Gamma \proves \gamma$, $\Delta \proves
\delta$ и $\gamma, \delta \proves \alpha$, то справедливо и $\Gamma, \Delta
\proves \alpha$.
\end{lmm}
\begin{proof}
Упражнение.
\end{proof}

Возьмем некоторую связку исчисления высказываний, например конъюнкцию:
$A\mathbin{\&}B$. Построим для нее таблицу истинности. По каждой строчке построим
утверждение, в котором отрицания появляются там, где в таблице истинности
находится $\False$:
\[
\begin{array}{cc|c|l}
A & B & A\mathbin{\&}B & \text{утверждение}\\\hline
\False & \False & \False & \neg A, \neg B \proves \neg(A\mathbin{\&}B)\\
\False & \True & \False & \neg A, B \proves \neg(A\mathbin{\&}B)\\
\True & \False & \False & A, \neg B \proves \neg(A\mathbin{\&}B)\\
\True & \True & \True & A, B \proves A\mathbin{\&}B
\end{array}
\]

\begin{lmm}
Каждое из построенных по таблицам истинности утверждений доказуемо.
\end{lmm}
\begin{proof}
Упражнение.
\end{proof}

\begin{lmm}[Правило контрапозиции]
Каковы бы ни были формулы $\alpha$ и $\beta$, справедливо, что
$\proves (\alpha \to \beta) \to (\neg\beta \to \neg\alpha)$.
\end{lmm}
\begin{proof}
Сперва докажем, что $\alpha \to \beta, \neg\beta \proves \neg\alpha$.
\begin{align*}
&(1)\ (\alpha \to \beta) \to (\alpha \to \neg\beta) \to \neg\alpha &&\text{Сх. акс. 9}\\
&(2)\ \alpha \to \beta &&\text{Допущение}\\
&(3)\ (\alpha \to \neg\beta) \to \neg\alpha &&\text{M.P. 2,1}\\
&(4)\ \neg\beta \to (\alpha \to \neg\beta) &&\text{Сх. акс. 1}\\
&(5)\ \neg\beta &&\text{Допущение}\\
&(6)\ \alpha \to \neg\beta &&\text{M.P. 5,4}\\
&(7)\ \neg\alpha &&\text{M.P. 6,3}
\end{align*}
Тогда, применив 2 раза Теорему о дедукции, получим вывод требуемого
утверждения.
\end{proof}

\begin{lmm}[Правило исключенного третьего]
Какова бы ни была формула $\alpha$, $\proves \alpha \vee \neg\alpha$.
\end{lmm}
\begin{proof}
Доказательство проведем в 3 этапа.
\begin{enumerate}
\item Для начала покажем $\proves \neg(\alpha\vee\neg\alpha) \to \neg\alpha$:
  \begin{align*}
  &(1)\ \alpha \to \alpha\vee\neg\alpha &&\text{Сх. акс. 6}\\
  &(2)\dots(n+1)\ \gamma_1, \dots \gamma_{n-1}, (\alpha \to \alpha\vee\neg\alpha) \to (\neg(\alpha\vee\neg\alpha) \to \neg\alpha) &&\text{Д-во из леммы 5.4}\\
  &(n+2)\ \neg(\alpha\vee\neg\alpha) \to \neg\alpha &&\text{M.P. } 1, n+1
  \end{align*}
\item Затем докажем $\proves \neg(\alpha\vee\neg\alpha) \to \neg\neg\alpha$:
  \begin{align*}
  &(1)\ \neg\alpha \to \alpha\vee\neg\alpha &&\text{Сх. акс. 7}\\
  &(2)\dots(k+1)\ \delta_1, \dots \delta_{k-1}, (\neg\alpha \to \alpha\vee\neg\alpha) \to (\neg(\alpha\vee\neg\alpha) \to \neg\neg\alpha) &&\text{Д-во из леммы 5.4}\\
  &(k+2)\ \neg(\alpha\vee\neg\alpha) \to \neg\neg\alpha &&\text{M.P. } 1, k+1
  \end{align*}
\item Теперь докажем все вместе:
  \begin{align*}
  &(1)\ \neg(\alpha\vee\neg\alpha) \to \neg\alpha &&\text{по пункту 1}\\
  &(2)\ \neg(\alpha\vee\neg\alpha) \to \neg\neg\alpha &&\text{по пункту 2}\\
  &(3)\ (\neg(\alpha\vee\neg\alpha) \to \neg\alpha) \to (\neg(\alpha\vee\neg\alpha) \to \neg\neg\alpha) \to (\neg\neg(\alpha\vee\neg\alpha)) &&\text{Сх. акс. 9}\\
  &(4)\ (\neg(\alpha\vee\neg\alpha) \to \neg\neg\alpha) \to \neg\neg(\alpha\vee\neg\alpha) &&\text{M.P. 1,3}\\
  &(5)\ \neg\neg(\alpha\vee\neg\alpha) &&\text{M.P. 2,4}\\
  &(6)\ \neg\neg(\alpha\vee\neg\alpha) \to (\alpha\vee\neg\alpha) &&\text{Сх. акс. 10}\\
  &(7)\ \alpha\vee\neg\alpha &&\text{M.P. 5,6}
  \end{align*}
\end{enumerate}
\end{proof}

\begin{lmm}[Об исключении допущения]
Пусть справедливо $\Gamma, \rho \proves \alpha$ и $\Gamma, \neg\rho \proves
\alpha$. Тогда также справедливо $\Gamma \proves \alpha$.
\end{lmm}
\begin{proof}
Упражнение.
\end{proof}

\begin{thm}[О полноте исчисления высказываний]
Пусть справедливо $\models \alpha$. Тогда также справедливо, что $\proves
\alpha$.
\end{thm}
\begin{proof}
Для доказательства теоремы мы докажем чуть более сильное утверждение~--- что
для любого $k$ от $0$ до $n$ и любой оценки переменных $x_1, \dots x_k$
справедливо $[x_1]P_1, \dots [x_k]P_k \proves \alpha$. Нетрудно заметить, что
утверждение теоремы непосредственно следует из данного утверждения для $k = 0$.
Доказательство будет вестись индукцией по $n - k$.

\emph{База.} Пусть $n - k = 0$, то есть $k = n$. $\models \alpha$ означает, что
при любой оценке $x_1, \dots x_n$ пропозициональных переменных $P_1, \dots P_n$
справедливо $\sem{\alpha}{P_1:=x_1,\dots P_n:=x_n} = \True$. Возьмем некоторую
оценку переменных $x_1, \dots x_n$. Тогда, по лемме 5,
$[x_1]P_1, \dots [x_n]P_n \proves [\sem{\alpha}{P_1:=x_1,\dots P_n:=x_n}]\,\alpha$,
то есть $[x_1]P_1, \dots [x_n]P_n \proves \alpha$.

\emph{Переход.} Пусть утверждение уже доказано для некоторого $n - k > 0$,
покажем его для $n - k + 1$ (то есть доказано для $k \le n$, покажем его для
$k - 1$). Возьмем некоторую оценку переменных $x_1, \dots x_{k-1}$. По
предположению, $[x_1]P_1, \dots [x_k]P_k \proves \alpha$, то есть
\begin{align*}
&[x_1]P_1, \dots [x_{k-1}]P_{k-1}, \neg P_k \proves \alpha\\
&[x_1]P_1, \dots [x_{k-1}]P_{k-1}, P_k \proves \alpha
\end{align*}
Тогда по лемме об исключении допущения, справедливо $[x_1]P_1, \dots
[x_{k-1}]P_{k-1} \proves \alpha$.
\end{proof}

\begin{thm}[О корректности исчисления высказываний]
Пусть справедливо $\proves \alpha$. Тогда также справедливо, что $\models
\alpha$.
\end{thm}
\begin{proof}
По условию теоремы, у нас есть доказательство высказывания $\alpha$, то есть
последовательность высказываний $\alpha_1, \dots \alpha_m$. Каждое
высказывание~--- это либо аксиома, либо применение правила Modus Ponens.
Докажем, что для каждого $k$ все высказывания $\alpha_l$ при $l \le k$~---
тавтологии. Доказательство будем вести индукцией по $k$.

\emph{База.} Пусть $k = 0$, тогда нет ни одного высказывания, про которое нужно
доказать, что оно~--- тавтология, то есть утверждение автоматически верно.

\emph{Переход.} Пусть для некоторого $k$ утверждение справедливо, докажем его
для $k+1$. Выберем некоторую оценку $x_1, \dots x_n$ пропозициональных
переменных $P_1, \dots P_n$, использованных в высказываниях $\alpha_1 \dots
\alpha_{k+1}$. Рассмотрим случаи.

Пусть $\alpha_{k+1}$~--- аксиома. В данную аксиому входят одна, две или три
формулы $\beta_1, \beta_2, \beta_3$. Подстановкой всех возможных истинностных
значений вместо данных формул можно проверить, что все аксиомы являются
тавтологиями, значит, они будут истинны и на тех конкретных значениях, которые
примут данные формулы после подстановки значений $x_1, \dots x_n$.

Пусть $\alpha_{k+1}$ получается по правилу Modus Ponens из $\alpha_p$ и
$\alpha_q$, причем $\alpha_q \equiv \alpha_p \to \alpha_{k+1}$. Тогда
$\sem{\alpha_p}{P_1:=x_1,\dots P_n:=x_n} = \True$ и $\sem{\alpha_p \to
\alpha_{k+1}}{P_1:=x_1,\dots P_n:=x_n} = \True$. Из таблицы истинности
импликации следует, что неизбежно $\sem{\alpha_{k+1}}{P_1:=x_1,\dots
P_n:=x_n} = \True$.
\end{proof}

Заметим, что вместе из этих двух теорем следует, что если неверно, что $\proves
\alpha$, то неизбежно найдется контрпример.

\section{Интуиционистское исчисление высказываний}
Одна из главных причин возникновения математической логики~--- кризис в
математике начала XX века. Классический пример такого~--- парадокс Рассела,
утверждающий, что понятие «множества всех не принадлежащих себе множеств»
противоречиво. В самом деле, пусть $X = \{t \mid t \notin t\}$. Рассмотрим $X
\in X$. Обе возможных альтернативы~--- $X \in X$ и $X \notin X$~--- вступают в
противоречие с определением множества $X$. Значит, мы можем сделать вывод,
множества $X$ не существует.

С одной стороны, в этом нет проблемы: мы не переживаем, если не найдём числа $t
\in \mathbb{N}$, такого, что $t = t$ и $t \ne t$ одновременно. Но, с другой
стороны, в формальном определении множества $X$ на первый взгляд нет никаких
проблем. Единственная, по-видимому, для современного читателя непривычная
деталь в данном определении~--- множество, принадлежащее самому себе, но если
вспомнить язык Java, то это возражение не должно вызывать никаких серьёзных
сомнений. Просто рассмотрите следующее определение:
\begin{verbatim}
class IntList {
   IntList next;
   int value;
}
\end{verbatim}

Парадокс Рассела был получен не сразу, а через несколько лет после появления
теории множеств и построения значительного количества математических теорий на
её основе. Из этого возникает сомнение, нет ли таких же противоречий и в
определении, например, вещественных чисел, просто глубже скрытых. Вдруг через
несколько лет какой-то математик найдёт противоречия в математическом анализе~---
и значительную часть теории придётся пересмотреть или вообще признать
ошибочной.

Было предложено много подходов к решению этой проблемы. Современный
(классический) подход состоит в том, чтобы так формализовать теорию множеств,
чтобы в ней не возникало парадоксов. Была сформулирована программа,
предполагавшая своей целью доказательство непротиворечивости такой формальной
теории. Однако, в 30-е годы Куртом Гёделем было показано, что такое
доказательство минимально удовлетворительной надёжности построить невозможно.
Конечно, самая распространённая формализация теории множеств~--- аксиоматика
Цермело--Френкеля~--- оформилась в современном виде в 1925 году, почти 100 лет
назад, и пока никаких противоречий в ней не было найдено. Но здесь всегда
остаётся место для сомнений.

Поэтому интерес представляют альтернативные подходы к проблеме. Данный раздел
посвящён интуиционистской логике. Математики, стоявшие у истоков данной логики,
видели решение в том, чтобы исключить из рассмотрения «неконструктивные»
объекты~--- объекты, метода построения которых не предложено. В частности,
множество $X$ из примера выше является примером неконструктивного объекта. Мы
слишком легко приняли на веру возможность его существования и получили
противоречие.

Столь резкие результаты получаются не всегда, но в целом в математике имеет
место довольно много утверждений, пусть и не противоречивых, но выглядящих
совершенно антиинтуитивно. Например, широко известна теорема Банаха--Тарского,
утверждающая, что трёхмерный шар можно разрезать на конечное число попарно
непересекающихся частей, из которых потом можно составить два шара того же
размера.

\subsection*{BHK-интерпретация}
BHK-интерпретация логики названа по именам математиков, её предложивших (Л.
Брауэр, А. Гейтинг и А. Колмогоров). Они решили изменить сам подход к
математическому рассуждению, предположив, что математик не думает в стиле
классической логики, и что правильно, поэтому, попробовать формализовать
«интуитивный» стиль.

Попробуем сформулировать эти соображения применительно к логическим связкам
исчисления высказываний. Будем определять интерпретацию индуктивно. Пусть даны
высказывания $P$ и $Q$, тогда:
\begin{itemize}
\item мы считаем $P\mathbin{\&}Q$ доказанным, если у нас есть доказательство
  высказывания $P$ и доказательство высказывания $Q$;
\item мы считаем $P\vee Q$ доказанным, если у нас есть доказательство $P$ или
  есть доказательство $Q$ (т.е. мы знаем, какая из двух альтернатив выполнена);
\item мы считаем $P \to Q$ доказанным, если мы умеем перестраивать любое
  доказательство высказывания $P$ в доказательство высказывания $Q$;
\item мы считаем $\bot$ утверждением, не имеющим доказательства.
\item $\neg P$ есть сокращение для $P \to \bot$. Иными словами, мы считаем
  $\neg P$ доказанным, если мы умеем из доказательства $P$ получить
  противоречие.
\end{itemize}

Проиллюстрировать подход можно на примере следующей теоремы.

\begin{thm}
Существуют два таких вещественных иррациональных числа $p$ и $q$, что $p^q \in
\mathbb{Q}$.
\end{thm}
\begin{proof}
Рассмотрим два случая:
\begin{enumerate}
\item если $\sqrt{2}^{\sqrt{2}} \in \mathbb{Q}$, то мы нашли требуемые числа:
  $p = q = \sqrt{2}$;
\item если же $\sqrt{2}^{\sqrt{2}} \notin \mathbb{Q}$, то рассмотрим
  $p = \sqrt{2}^{\sqrt{2}}$ и $q = \sqrt{2}$, тогда
  $p^q = \left(\sqrt{2}^{\sqrt{2}}\right)^{\sqrt{2}} = \sqrt{2}^2 = 2 \in
  \mathbb{Q}$.
\end{enumerate}
\end{proof}

Данная теорема, хоть и доказывает факт существования таких чисел, ничего не
говорит по поводу того, какой из двух случаев имеет место~--- то есть, она
неконструктивна. В самом деле, обозначим факт того, что $\sqrt{2}^{\sqrt{2}}
\in \mathbb{Q}$ за $P$, а итоговое утверждение теоремы~--- за $S$. Мы
показываем, что $P \to S$, и что $\neg P \to S$. Однако, чтобы перейти к просто
$S$, нам нужно показать $P \vee \neg P$. Несложно видеть, что в
BHK-интерпретации нет простого способа это сделать: чтобы считать дизъюнкцию
доказанной, мы должны знать, какой из случаев имеет место. Поэтому, данное
рассуждение не является доказательством в BHK-интерпретации.

Для программистов же здесь важным является следующее соображение: эта теорема
не позволяет написать программу, ищущую эти два числа. Скажем, теорема о
дедукции не такова: её доказательство позволяет построить такую программу,
предъявляющую объект, существование которого утверждает теорема.

\subsection*{Формализация интуиционистской логики}
Исходный постулат интуиционистской логики состоит в том, что никакая
формализация не является первичной. Мы выбираем те или иные правила только
потому, что они соответствуют заявленным целям. Мы также вольны в любой момент
правила поменять, если на то будут серьёзные основания.

У интуиционистской логики есть несколько формализаций, рассмотрим наиболее
распространённую. Заменим аксиому устранения двойного отрицания (10-ю) на
$\alpha \to \neg\alpha \to \beta$. Полученную систему назовём аксиоматикой
интуиционистского исчисления высказываний.

А у данной аксиоматики есть интересные свойства, отсутствующие в классическом
исчислении высказываний. Например, если $\proves_{\text{И}} \alpha \vee \beta$,
то $\proves_{\text{И}} \alpha$, либо $\proves_{\text{И}} \beta$~--- сравните с
BHK-интерпретацией дизъюнкции. Данное следствие поясняет обоснованность замены
аксиомы, в дальнейшем оно будет доказано формально.

\subsection{Булева алгебра и топологическая интерпретация интуиционистского исчисления высказываний}
Мы построим две параллельные интерпретации для классической и
интуиционистской логики.

\begin{dfn}
Пусть дано некоторое исчисление высказываний, для которого нам нужно построить
модель~--- предложить способ оценки истинности выражений. Начинаем мы с
множества истинностных значений. Возьмем в качестве этого множества все
открытые множества некоторого заранее выбранного топологического пространства.
Определим оценку для связок интуиционистского исчисления высказываний следующим
образом:
\begin{align*}
\llbracket A\mathbin{\&}B\rrbracket &= \llbracket A\rrbracket \cap \llbracket B\rrbracket\\
\llbracket A \vee B\rrbracket &= \llbracket A\rrbracket \cup \llbracket B\rrbracket\\
\llbracket A \to B\rrbracket &= (\cmpl\llbracket A\rrbracket \cup \llbracket B\rrbracket)^\circ\\
\llbracket \neg A\rrbracket &= (\cmpl\llbracket A\rrbracket)^\circ
\end{align*}
Будем считать, что формула истинна в данной модели, если её значение оказалось
равно всему пространству.
\end{dfn}

Например, возьмем в качестве пространства $\mathbb{R}$, и вычислим значение
формулы $A \vee \neg A$ при $A$ равном $(0, 1)$:
\begin{multline*}
\llbracket A\vee\neg A\rrbracket = (0,1)\cup\llbracket\neg A\rrbracket = (0,1)\cup(\cmpl(0,1))^\circ\\
= (0,1)\cup((-\infty,0)\cup(1,\infty)) = (-\infty,0)\cup(0,1)\cup(1,\infty).
\end{multline*}
Нетрудно видеть, что данная формула оказалась не общезначимой в данной
интерпретации.

\begin{thebibliography}{1}
\bibitem{viro} Виро О.Я., Иванов А.О., Нецветаев Н.Ю., Харламов В.М.
  Элементарная топология --- М.: МЦНМО, 2012
\end{thebibliography}

\end{document}
